\documentclass[11pt]{article}
\usepackage[margin=0.95in]{geometry}
\usepackage{amsmath,amssymb}
\usepackage{graphicx}
\usepackage{booktabs,longtable,array}
\usepackage{xcolor}
\usepackage[hidelinks]{hyperref}
\usepackage{enumitem}
\usepackage{caption}
\setlist{itemsep=2pt,topsep=4pt}
\captionsetup{font=small,labelfont=bf}
\setlength{\parskip}{5pt}
\setlength{\parindent}{0pt}
\renewcommand{\arraystretch}{1.15}
\title{Phase 5 Summary: Does the Actor Learning Rate Explain TRM's TD3 Failure?}
\author{}
\date{trm-delay-robust-scaling $\cdot$ phase summary $\cdot$ 2026-10-10 $\cdot$ 36 Katana jobs (+ 12 reused Stage T runs)\\[2pt]Labels are defined in \texttt{.\allowbreak{}.\allowbreak{}/\allowbreak{}glossary.\allowbreak{}md}. Pre-registration: \texttt{.\allowbreak{}.\allowbreak{}/\allowbreak{}.\allowbreak{}.\allowbreak{}/\allowbreak{}plan/\allowbreak{}plan.\allowbreak{}md}; rules and amendment: \texttt{.\allowbreak{}.\allowbreak{}/\allowbreak{}.\allowbreak{}.\allowbreak{}/\allowbreak{}.\allowbreak{}.\allowbreak{}/\allowbreak{}experiments/\allowbreak{}phase5/\allowbreak{}gates.\allowbreak{}json}; analysis: \texttt{.\allowbreak{}.\allowbreak{}/\allowbreak{}.\allowbreak{}.\allowbreak{}/\allowbreak{}.\allowbreak{}.\allowbreak{}/\allowbreak{}experiments/\allowbreak{}analysis\_\allowbreak{}p5\_\allowbreak{}t2.\allowbreak{}json}.}
\begin{document}
\maketitle
\section*{Question}



Phase 4's Stage T reran Phase 3's capacity contrast under TD3 with its standard settings (\texttt{td3.yaml}). GT passed, TRM 788k $-$ MLP 788k negative in 12/12, but by an optimisation failure:

\begin{itemize}
\item TRM 788k stayed at random-policy level in 10/12 runs;
\item its critic loss rose through training;
\item both MLPs learned in every run.
\end{itemize}



The suspected cause was the actor learning rate. \texttt{td3.yaml} uses 1e-3, 10$\times$ the DDPG actor rate (1e-4) under which TRM trained normally in Phases 1--4, and Phase 4's LR sweep showed that a higher actor LR hurts TRM. Phase 5 tests that one explanation.



\section*{Methodology}



\begin{itemize}
\item \textbf{Pre-registration.} The plan and gate GT2 were written into \texttt{plan.md} and \texttt{gates.json} before any data, under the same procedure as Phase 4 (frozen code snapshot \texttt{p5code/t2}, its own ledger, stop at the gate). Amendment P5-1, made before data with the author's approval, reuses Stage T's 12 TRM runs as the actor-LR 1e-3 arm.
\item \textbf{Code.} No changes. \texttt{src/} and \texttt{conf/} are byte-identical to Phase 4's frozen Stage T snapshot.
\item \textbf{Setting.} Pendulum-v1, TD3 (\texttt{td3.yaml}), with only \texttt{arch.actor\_lr} overridden; the critic LR stays 1e-3. 200 episodes, conditions C0 (no delay) and C1 (fixed d = 1 + queue, no previous action), seeds 1--6.
\end{itemize}



\begin{center}\begin{tabular}{llrl}
\toprule
Arm & Actor LR & Runs & Source \\
\midrule
TRM 788k & 1e-4 & 12 & new \\
TRM 788k & 3e-4 & 12 & new \\
TRM 788k & 1e-3 & 12 & Stage T (P5-1) \\
MLP 788k (control) & 1e-4 & 12 & new \\
\bottomrule
\end{tabular}\end{center}



\begin{itemize}
\item \textbf{Rule (GT2).}
\begin{itemize}
\item "Learned" = final avg-100 > $-$1000.
\item "LR explains the TD3 failure" if TRM 788k learns in $\geq$ 10/12 runs at some actor LR below 1e-3.
\item If so, TRM 788k $-$ MLP 788k at that LR is reported (n = 12 pairs, sign test). The LR-matched 1e-4 contrast is primary.
\end{itemize}
\end{itemize}



\section*{Results}



\begin{center}\small\begin{tabular}{>{\raggedright\arraybackslash}p{0.316\linewidth}>{\raggedleft\arraybackslash}p{0.081\linewidth}>{\raggedleft\arraybackslash}p{0.193\linewidth}>{\raggedleft\arraybackslash}p{0.154\linewidth}>{\raggedleft\arraybackslash}p{0.107\linewidth}}
\toprule
Arm & Learned & Median final avg-100 & Late critic loss & Eval last-5 \\
\midrule
TRM 788k, actor LR 1e-4 & \textbf{12/12} & $-$187.7 & 36.0 & $-$173.1 \\
TRM 788k, actor LR 3e-4 & 11/12 & $-$234.6 & 47.4 & $-$200.7 \\
TRM 788k, actor LR 1e-3 (Stage T) & 2/12 & $-$1389.3 & 560.9 & $-$1415.1 \\
MLP 788k, actor LR 1e-4 & 12/12 & $-$176.1 & 25.4 & $-$168.0 \\
MLP 788k, actor LR 1e-3 (Stage T) & 12/12 & $-$175.9 & 23.7 & $-$168.3 \\
\bottomrule
\end{tabular}\end{center}



\begin{center}\small\begin{tabular}{>{\raggedright\arraybackslash}p{0.523\linewidth}>{\raggedleft\arraybackslash}p{0.069\linewidth}>{\raggedleft\arraybackslash}p{0.160\linewidth}>{\raggedleft\arraybackslash}p{0.126\linewidth}}
\toprule
Contrast (n = 12) & Median & Pairs positive & Two-sided p \\
\midrule
\textbf{TRM 788k $-$ MLP 788k, both at actor LR 1e-4 (primary)} & \textbf{$-$13.0} & \textbf{2/12} & \textbf{0.039} \\
TRM 788k at 3e-4 $-$ MLP 788k at 1e-4 (unequal LRs) & $-$53.6 & 0/12 & 0.0005 \\
\bottomrule
\end{tabular}\end{center}



\textbf{Gate GT2: "LR explains".} TRM 788k learns in 12/12 runs at actor LR 1e-4, against 2/12 at TD3's default of 1e-3.



\section*{Discussion}



\begin{itemize}
\item \textbf{Stage T's failure was an actor-learning-rate effect, not something specific to TD3.}
\begin{itemize}
\item At the DDPG-like actor LR, TRM trains under TD3 in every run.
\item At 3e-4 it trains in all but one run (C0 seed 1 stays near random).
\item At 1e-3 it mostly fails.
\item The critic-loss trend matches: 36 $\rightarrow$ 47 $\rightarrow$ 561 as the actor LR rises.
\end{itemize}
\item \textbf{With the LR fixed, Phase 3's attribution replicates under a second algorithm.} TRM 788k $-$ MLP 788k is $-$13.0 under TD3 (2/12 pairs, p = 0.039). Phase 3's G1 under DDPG gave $-$27.1 (2/12). The direction and sign count are the same; the median is about half the size, with almost the same seed spread. TRM's value fit is again worse (median late critic loss 36.0 vs 25.4, about 1.4$\times$; the median per-pair ratio is 1.5$\times$, in line with Phase 3's 1.45--1.89$\times$).
\item \textbf{TRM is far more sensitive to the actor learning rate than the MLP.} The MLP is unaffected by a 10$\times$ change ($-$176.1 at 1e-4, $-$175.9 at 1e-3). TRM goes from 12/12 learned to 2/12. This is a practical cost of the recursive trunk: a standard algorithm configuration that suits an MLP can break it.
\item \textbf{This resolves Phase 4's GT caveat.} At matched, sensible learning rates, TRM shows the same modest, consistent deficit under DDPG and TD3, rather than an algorithm-specific failure.
\end{itemize}



\section*{Threats}



\begin{itemize}
\item The "learned" cut ($-$1000) is coarse. It separates learning from staying near random ($\approx$ $-$1164), not good from mediocre runs.
\item Three learning rates only. 1e-4 is not shown to be TRM's best TD3 setting.
\item The MLP control runs at 1e-4 only, so the 3e-4 contrast mixes architecture with LR and is secondary.
\item The 1e-3 arm is reused from Stage T. The configuration and code are identical, but runs are not bit-reproducible across GPU models.
\end{itemize}



\section*{What Phase 5 changes}



\begin{center}\small\begin{tabular}{>{\raggedright\arraybackslash}p{0.467\linewidth}>{\raggedright\arraybackslash}p{0.467\linewidth}}
\toprule
Before Phase 5 & After Phase 5 \\
\midrule
G1's direction holds under TD3, but via an optimisation failure (GT caveat) & G1 replicates under TD3 at a matched actor LR: $-$13.0, 2/12, p = 0.039 \\
Likely cause of the TD3 failure (actor LR) untested & Tested and confirmed (GT2) \\
--- & New: TRM is markedly more sensitive to the actor learning rate than the MLP \\
\bottomrule
\end{tabular}\end{center}



\textbf{Bottom line:} under both DDPG and TD3, with sensible learning rates, TRM 788k trails a parameter-matched MLP by a small, consistent margin and fits its value function worse. It is also much more fragile to the actor learning rate.
\end{document}
